% ==============================================================================
% METADATOS DE AUDITORÍA EDITORIAL Y RIGOR CIENTÍFICO (PROTOCOLO AGENTS.md)
% ==============================================================================
% Capítulo: 02
% Título: Alfabetización Digital e Inteligencia Artificial en la Educación Superior: Fundamentos, Marcos de Competencia y Mitigación Epistémica
% Grupo de Origen: Grupo 12 (Asignatura de Inteligencia Artificial Aplicada, UNMSM)
% Archivo Fuente Analizado: Capitulo_2_Alfabetizacion_digital_final.pdf (35 páginas)
% Fuentes Primarias Cotejadas y Verificadas (referencias/):
%   1. Cosgrove, J., & Cachia, R. (2025). DigComp 3.0: European Digital Competence Framework (Publications Office of the European Union, DOI: 10.2760/0001149)
%   2. Vuorikari, R., Kluzer, S., & Punie, Y. (2022). DigComp 2.2: The Digital Competence Framework for Citizens (Publications Office of the European Union, DOI: 10.2760/115376)
%   3. Carolus, A., Koch, M. J., Straka, S., Latoschik, M. E., & Wienrich, C. (2023). MAILS – Meta AI Literacy Scale (Computers in Human Behavior: Artificial Humans, DOI: 10.1016/j.chbah.2023.100014)
%   4. Ng, D. T. K., Leung, J. K. L., Chu, S. K. W., & Qiao, M. S. (2021). Conceptualizing AI literacy: An exploratory review (Computers and Education: AI, DOI: 10.1016/j.caeai.2021.100041)
%   5. Pennycook, G., & Rand, D. G. (2022). Accuracy prompts are a replicable and generalizable approach for reducing the spread of misinformation (Nature Communications, DOI: 10.1038/s41467-022-30073-5)
%   6. Roozenbeek, J., van der Linden, S., Goldberg, B., Rathje, S., & Lewandowsky, S. (2022). Psychological inoculation improves resilience against misinformation (Science Advances, DOI: 10.1126/sciadv.abo6254)
%   7. Temoshok, D. et al. (2025). Digital Identity Guidelines: Authentication and Authenticator Management (NIST SP 800-63B-4, DOI: 10.6028/NIST.SP.800-63B-4)
%   8. Ara, J., Sik-Lanyi, C., & Kelemen, A. (2024). Accessibility engineering in web evaluation process (Universal Access in the Information Society, DOI: 10.1007/s10209-023-00967-2)
%   9. Tinmaz, H., Lee, Y.-T., Fanea-Ivanovici, M., & Baber, H. (2022). A systematic review on digital literacy (Smart Learning Environments, DOI: 10.1186/s40561-022-00204-y)
%  10. Khan, S., Webster, S., Puxty, J., & Robertson, M. (2026). Key Challenges and Barriers to Digital Literacy for Older Adults (JMIR Aging, DOI: 10.2196/80647)
% Cláusula NotebookLM: Grounding estricto al 100%. Cero datos o métricas inventadas.
% Parámetro de Extensión: >= 25 páginas compiladas en formato oficial A4 UNMSM.
% Decisiones del Usuario: Eliminación total de la columna lateral persistente de definiciones e integración holística en cuerpo continuo, tablas formalizadas y cajas institucionales.
% Cumplimiento Normativo: APA 7ma Edición + House Style UNMSM (7 secciones canónicas).
% Estado de Auditoría: CONFORME / REVISADO POR CAPITULO-AUDITOR
% Fecha de Auditoría: 2026-11-20
% ==============================================================================

\section[Alfabetización Digital e IA en Educación]{Alfabetización Digital e Inteligencia Artificial en la Educación Superior: Fundamentos, Marcos de Competencia y Mitigación Epistémica}
\markboth{CAPÍTULO 2\quad ALFABETIZACIÓN DIGITAL E IA EN EDUCACIÓN}{}

\begin{cajabeige}
\textbf{Pregunta Central del Capítulo:} \textquestiondown Cómo deben articularse los marcos globales de competencia digital ciudadana y docente (DigComp 3.0, DigCompEdu y UNESCO) con la alfabetización crítica en inteligencia artificial (AI Literacy) para transformar las prácticas pedagógicas en la educación superior, garantizando al mismo tiempo la accesibilidad universal, la soberanía de los datos académicos y la resiliencia epistémica frente a la desinformación sintética?\\[0.3cm]
\textbf{Objetivo de Aprendizaje:} El lector analizará la evolución epistemológica de la alfabetización digital, examinará la estructura multidimensional de los marcos DigComp 2.2, 3.0 y DigCompEdu, evaluará intervenciones de inoculación psicológica y pausas de precisión contra la desinformación, aplicará directrices criptográficas NIST SP 800-63B-4 y estándares WCAG 2.2 de accesibilidad, y medirá el perfil de competencias en IA mediante la escala psicométrica MAILS en contextos universitarios.
\end{cajabeige}

\textoside{
La alfabetización digital e inteligencia artificial en la sociedad contemporánea constituye una praxis sociotécnica indivisible: amalgama la destreza operacional para interactuar con sistemas algorítmicos complejos con una capacidad axiológica de juicio crítico para evaluar la veracidad, procedencia y sesgos de la información. No se reduce al dominio técnico instrumental de interfaces o suites ofimáticas, sino que fundamenta la agencia epistémica ciudadana, la autonomía intelectual y la preservación de los derechos fundamentales en ecosistemas de información mediados por modelos predictivos estocásticos.
}{CONCEPTO CLAVE\\[0.1cm]Alfabetización Sociotécnica y Agencia Epistémica}

\vspace{0.4cm}

\subsection{Epistemología y Fundamentos de la Alfabetización Digital en la Sociedad Algorítmica}

\subsubsection{Evolución Conceptual: De la Destreza Instrumental a la Agencia Sociotécnica}

La comprensión teórica de la alfabetización en entornos computacionales ha experimentado una profunda metamorfosis conceptual a lo largo de las últimas cinco décadas. En las etapas fundacionales de la informática personal durante las décadas de 1970 y 1980, autores seminales como Arthur Luehrmann (1972) y Seymour Papert (1980) concibieron la \textbf{alfabetización computacional} (\textit{computer literacy}) como la capacidad de programar, interactuar con líneas de comandos y comprender la arquitectura interna de los procesadores. Papert sostenía que programar en lenguajes como LOGO transformaba al ordenador en un artefacto cognitivo para pensar sobre el propio pensamiento, superando la visión pasiva de la máquina como mera instructora asistida.

Sin embargo, la emergencia y masificación de la World Wide Web a mediados de la década de 1990 reconfiguró radicalmente las demandas formativas de la ciudadanía. En su obra seminal, Paul Gilster (1997) acuñó formalmente el término \textbf{alfabetización digital} (\textit{digital literacy}), definiéndola primordialmente como la habilidad para comprender y utilizar información proveniente de una multiplicidad de fuentes presentadas a través de ordenadores, enfatizando que el aspecto central no radicaba en las destrezas mecánicas del teclado o el ratón, sino en el dominio del pensamiento crítico aplicado al contenido hipertextual.

Como evidencian Tinmaz, Lee, Fanea-Ivanovici y Baber (2022) en su exhaustiva revisión sistemática de la literatura, el constructo contemporáneo de la alfabetización digital trasciende con creces la dicotomía tradicional entre usuarios técnicos y no técnicos. Los autores demuestran que una competencia digital madura opera como una estructura tripartita e indisociable:
\begin{itemize}
    \item \textbf{Dimensión Cognitiva (Conocimiento / \textit{Saber}):} Comprensión conceptual de los principios de funcionamiento de las redes, los protocolos de comunicación, la arquitectura de los datos y las dinámicas socioeconómicas del capitalismo de plataformas.
    \item \textbf{Dimensión Instrumental y Procedimental (Habilidad / \textit{Saber Hacer}):} Habilidad técnica y fluidez operacional para navegar, buscar, crear, filtrar, proteger y configurar dispositivos, plataformas y flujos de trabajo computacionales.
    \item \textbf{Dimensión Axiológica y Crítica (Actitud y Juicio / \textit{Saber Juzgar}):} Disposición ética y reflexiva para sopesar la veracidad de los enunciados, evaluar las intenciones del emisor, anticipar los riesgos de privacidad, regular el bienestar socioemocional y participar responsablemente en la esfera pública democrática.
\end{itemize}

Esta visión holística es respaldada por Valverde-Berrocoso, González-Fernández y Acevedo-Borrega (2022), quienes señalan que en los ecosistemas infocomunicativos hiperconectados contemporáneos resulta indispensable hablar de \textbf{multialfabetización}. La confluencia de la alfabetización informacional, mediática, visual y algorítmica impide que la mera familiaridad superficial con dispositivos táctiles sea confundida con una auténtica madurez epistémica.

\subsubsection{Acceso Material, Mediación Instrumental y Capacidad de Juicio Crítico}

El despliegue de cualquier actividad en el ciberespacio depende de una interacción jerárquica entre condiciones materiales de acceso, habilidades procedimentales de manejo y marcos analíticos de discernimiento crítico. La literatura empírica demuestra de forma contundente que poseer acceso material no garantiza el desarrollo de habilidades evaluativas. 

Para desentrañar esta interacción, es imperativo establecer distinciones técnicas rigurosas entre los componentes que articulan la mediación instrumental cotidiana:
\begin{enumerate}
    \item \textbf{Navegador Web (\textit{Web Browser}):} Aplicación de software cliente (como Chromium, Firefox o WebKit) encargada de interpretar código HTML, estilos CSS y scripts JavaScript, gestionar cookies de sesión, aislar pestañas mediante cajas de arena criptográficas (\textit{sandboxing}) y procesar solicitudes seguras vía HTTPS.
    \item \textbf{Motor de Búsqueda (\textit{Search Engine}):} Sistema informático distribuido masivo (como Google Search o Bing) que rastrea la web mediante arañas algorítmicas (\textit{web crawlers}), indexa billones de páginas mediante vectores de palabras clave y clasifica los resultados a través de algoritmos deterministas y estocásticos (como PageRank y transformadores neuronales).
    \item \textbf{Localizador Uniforme de Recursos (\textit{URL}):} Cadena de caracteres estandarizada (RFC 3986) que identifica de manera unívoca un recurso en la red mediante protocolo (\texttt{https://}), nombre de dominio canónico (\texttt{unmsm.edu.pe}), puerto opcional y ruta jerárquica de acceso.
    \item \textbf{Indexador Académico Especializado:} Repositorios y bases de datos bibliométricas (como Scopus, Web of Science, SciELO o PubMed) que estructuran la literatura científica mediante metadatos normalizados (autoría, filiación institucional, resumen, palabras clave y citaciones cruzadas) filtrados por arbitraje por pares ciegos.
\end{enumerate}

La confusión generalizada entre estas herramientas introduce distorsiones epistémicas severas. Una persona que confunde la barra de navegación del navegador con la caja de entrada de un motor de búsqueda queda expuesta a la intermediación publicitaria. Como documenta la Organización para la Cooperación y el Desarrollo Económicos (OCDE / OECD, 2022), los motores de búsqueda comerciales operan bajo modelos económicos basados en la monetización de la atención, donde los primeros resultados visibles corresponden frecuentemente a enlaces patrocinados y optimización SEO agresiva, desplazando portales institucionales legítimos y repositorios académicos sin fines de lucro. El criterio crítico del usuario actúa aquí como la primera línea de defensa para discernir la intencionalidad algorítmica de la plataforma.

\subsubsection{Brecha Digital Multidimensional, Inclusión y Aprendizaje a lo Largo de la Vida}

Durante décadas, las políticas públicas gubernamentales concibieron la \textbf{brecha digital} como un problema estrictamente infraestructural, postulando erróneamente que la mera dotación masiva de computadoras o conexiones de fibra óptica cerraría automáticamente las asimetrías de conocimiento. La investigación contemporánea desmiente de forma categórica este tecno-optimismo reduccionista. En la actualidad, la comunidad científica reconoce un modelo de tres niveles jerárquicos de brecha digital:
\begin{itemize}
    \item \textbf{Brecha de Primer Nivel (Acceso Físico e Infraestructura):} Desigualdad material en la disponibilidad de hardware adecuado, ancho de banda simétrico, estabilidad del suministro eléctrico y acceso a terminales no compartidos.
    \item \textbf{Brecha de Segundo Nivel (Competencias y Calidad de Uso):} Asimetría en las habilidades cognitivas y operacionales requeridas para explotar las herramientas digitales. Incluye la capacidad para formular búsquedas complejas, verificar fuentes, automatizar tareas y resolver fallas técnicas elementales.
    \item \textbf{Brecha de Tercer Nivel (Capitalización de Resultados y Retorno Socioeconómico):} Diferencia sustancial en la capacidad de los individuos para traducir el uso de las tecnologías digitales en beneficios tangibles para el bienestar personal, el avance educativo, el empleo calificado y la participación cívica efectiva.
\end{itemize}

Un sector poblacional que experimenta con particular agudeza estas fricciones es el de los adultos mayores y jubilados. En un estudio empírico seminal publicado en \textit{JMIR Aging}, Khan, Webster, Puxty y Robertson (2026) sintetizan los desafíos y barreras que condicionan la alfabetización digital en etapas avanzadas de la vida. Los autores demuestran que las dificultades observadas no se deben a una supuesta aversión intrínseca a la innovación, sino a la conjunción de barreras cognitivas (sobrecarga por interfaces dinámicas no lineales), barreras sensoriales y motoras (disminución de la agudeza visual, pérdida de destreza motriz fina para pantallas táctiles diminutas), barreras afectivas (ansiedad informática, temor a cometer errores irreversibles o sufrir estafas financieras) y, fundamentalmente, a la escasez de redes de apoyo social y soporte técnico paciente.

Frente a esta realidad, Khan et al. (2026) postulan el principio pedagógico del \textbf{acompañamiento sin sustitución de decisiones}. Cuando un familiar o tutor resuelve un trámite digital quitándole el dispositivo de las manos a la persona mayor, priva al aprendiz de la experiencia cenestésica y reflexiva necesaria para consolidar el esquema cognitivo, profundizando la dependencia técnica. La verdadera inclusión digital exige andamiajes formativos que respeten los ritmos individuales, permitan el ensayo y error en entornos controlados y enfaticen la autonomía decisional.

Asimismo, la evidencia acumulada derriba de forma concluyente el extendido mito de los "nativos digitales" propuesto a inicios de siglo por Prensky (2001). Como subrayan Tinmaz et al. (2022) y Valverde-Berrocoso et al. (2022), el hecho de que las generaciones jóvenes hayan nacido inmersas en entornos con teléfonos inteligentes y redes sociales no les confiere espontáneamente capacidades para evaluar la robustez metodológica de un artículo científico, identificar una falacia argumentativa en un foro web o configurar directivas criptográficas de seguridad. La competencia digital no es una propiedad biológica generacional, sino una construcción pedagógica que debe ser explícita y rigurosamente cultivada en la educación superior.

\newpage

\subsection{Marcos Globales de Competencia Digital Docente y Ciudadana: De DigComp 2.2/3.0 a DigCompEdu y UNESCO}

\subsubsection{Anatomía de los Marcos Europeos DigComp 2.2 y DigComp 3.0}

Para estandarizar, evaluar y orientar el desarrollo de las habilidades tecnológicas en la ciudadanía, la Comisión Europea formuló el Marco de Competencia Digital para la Ciudadanía (\textit{European Digital Competence Framework for Citizens}, conocido globalmente como \textbf{DigComp}). Desde su versión preliminar en 2013, DigComp se ha consolidado como el estándar de referencia internacional más influyente, sirviendo de base para políticas curriculares, marcos de certificación laboral y diseño de planes de estudio universitarios en más de cincuenta países.

El marco se estructura de manera canónica en cinco grandes áreas de competencia que abarcan veintiuna subcompetencias operativas:
\begin{enumerate}
    \item \textbf{Área 1: Búsqueda, Filtrado y Gestión de Información y Datos:} Articula la navegación estratégica en la red, la formulación de consultas avanzadas, la evaluación crítica de la credibilidad de fuentes y contenidos, y el almacenamiento, organización y recuperación estructurada de datos digitales.
    \item \textbf{Área 2: Comunicación y Colaboración en Entornos Digitales:} Comprende la interacción mediante tecnologías diversas, el intercambio oportuno de recursos, la participación ciudadana en línea, la colaboración mediante herramientas compartidas, el cumplimiento de normas de convivencia digital (\textit{netiqueta}) y la gestión soberana de la identidad digital propia.
    \item \textbf{Área 3: Creación de Contenidos Digitales:} Abarca el desarrollo y edición de artefactos multimediales en múltiples formatos, la integración y reelaboración de contenidos preexistentes, la aplicación rigurosa de licencias de propiedad intelectual y derechos de autor (Creative Commons, copyright), y la formulación de instrucciones algorítmicas elementales (programación y flujos lógicos).
    \item \textbf{Área 4: Seguridad y Protección de Activos:} Involucra la protección física y lógica de dispositivos electrónicos, la preservación de datos personales y privacidad, la salvaguarda de la salud y el bienestar psicoemocional (prevención de ciberacoso y tecnoestrés), y la comprensión del impacto ambiental de la tecnología (huella de carbono de centros de datos y basura electrónica).
    \item \textbf{Área 5: Resolución de Problemas:} Desarrolla la resolución autónoma de fallos técnicos de hardware y software, la identificación analítica de necesidades de actualización tecnológica, el uso creativo y productivo de herramientas para innovar procesos, y la identificación de vacíos en la propia competencia digital para emprender aprendizajes autodirigidos.
\end{enumerate}

En la actualización DigComp 2.2, dirigida por Vuorikari, Kluzer y Punie (2022), el marco incorporó más de doscientos cincuenta ejemplos descriptivos de conocimientos, habilidades y actitudes, articulados en ocho niveles de dominio progresivo (desde el nivel inicial A1 hasta el nivel altamente especializado C2). Esta taxonomía permite a las instituciones educativas no solo certificar lo que un estudiante "sabe hacer", sino también evaluar su grado de autonomía e independencia al resolver tareas con diferente nivel de complejidad contextual.

Por su parte, la reciente publicación del marco \textbf{DigComp 3.0}, sintetizado por Cosgrove y Cachia (2025), representa una inflexión metodológica trascendental. A diferencia de las ediciones previas, donde la inteligencia artificial se abordaba como un tópico marginal de resolución técnica, DigComp 3.0 transversaliza la competencia en IA a lo largo de las cinco áreas de competencia. En el nuevo marco, se explicita que un ciudadano digital competente debe comprender la naturaleza predictiva y generativa de los modelos de lenguaje masivo (LLMs), evaluar la opacidad de los algoritmos de recomendación, identificar imágenes y audios sintéticos ultrafalsos (\textit{deepfakes}), gestionar la propiedad intelectual de obras co-creadas con asistentes artificiales y auditar la discriminación algorítmica en procesos de reclutamiento laboral o evaluación académica.

En el ámbito macroestadístico, la Unión Europea introdujo el Indicador de Habilidades Digitales 2.0 (\textit{Digital Skills Indicator 2.0}, DSI 2.0), documentado por Vuorikari, Jerzak, Karpinski, Pokropek y Tudek (2022). A partir de encuestas poblacionales armonizadas por Eurostat, el DSI 2.0 cuantifica el porcentaje de población que posee competencias digitales básicas o avanzadas en cada una de las dimensiones de DigComp. La investigación subraya la necesidad metodológica de no extrapolar mecánicamente indicadores macroestadísticos agregados a evaluaciones de desempeño individual: mientras el DSI 2.0 mide la prevalencia de uso reportado de funciones tecnológicas en la sociedad, la evaluación formativa universitaria exige instrumentos psicométricos de observación directa y resolución de problemas prácticos en el aula.

\subsubsection{El Marco de Competencia Digital Docente (DigCompEdu) y UNESCO ICT-CFT}

En el contexto específico de la educación superior, las competencias digitales ciudadanas resultan insuficientes para abordar las exigencias pedagógicas contemporáneas. Los docentes universitarios no solo deben ser usuarios competentes de tecnologías, sino también ingenieros del aprendizaje capaces de diseñar experiencias didácticas enriquecidas digitalmente. Para responder a este desafío, el Centro Común de Investigación (JRC) de la Comisión Europea desarrolló el Marco de Competencia Digital Docente (\textbf{DigCompEdu}), formulado por Redecker (2017).

DigCompEdu organiza la labor pedagógica en seis áreas diferenciadas y complementarias:
\begin{itemize}
    \item \textbf{Área 1: Compromiso Profesional:} Utilización de tecnologías digitales para la comunicación institucional reflexiva, la colaboración interdisciplinaria entre docentes, el desarrollo profesional continuo y la investigación académica colaborativa.
    \item \textbf{Área 2: Contenidos Digitales:} Selección crítica, curaduría fundamentada, creación, modificación, licenciamiento y distribución segura de recursos educativos digitales accesibles.
    \item \textbf{Área 3: Enseñanza y Aprendizaje:} Diseño, implementación y mediación de intervenciones didácticas digitales, promoviendo el aprendizaje colaborativo, la autorregulación del estudiante y la personalización de itinerarios formativos.
    \item \textbf{Área 4: Evaluación Digital:} Empleo de instrumentos y analíticas de aprendizaje para la evaluación formativa y sumativa, garantizando la retroalimentación oportuna, la validez probatoria y la equidad en la calificación.
    \item \textbf{Área 5: Empoderamiento de los Estudiantes:} Aplicación de principios de Diseño Universal para el Aprendizaje (DUA) para mitigar barreras de accesibilidad, responder a la neurodiversidad, atender necesidades educativas específicas y fomentar el compromiso activo del alumnado.
    \item \textbf{Área 6: Facilitación de la Competencia Digital de los Estudiantes:} Acciones pedagógicas explícitas orientadas a que los estudiantes desarrollen su propia alfabetización informacional, comunicación en red, creación multimedia, ciberseguridad y resolución crítica de problemas.
\end{itemize}

De manera complementaria, el Marco de Competencias de los Docentes en Materia de TIC de la UNESCO (ICT-CFT) y las recientes \textit{Directrices de la UNESCO para la IA y la Educación} (2024) refuerzan este paradigma insistiendo en el principio de la \textbf{pedagogía centrada en lo humano}. La UNESCO advierte categóricamente contra la automatización indiscriminada del acto educativo: los docentes deben salvaguardar la relación intersubjetiva, utilizar los modelos generativos como andamiajes auxiliares y preservar espacios de debate deliberativo donde la argumentación racional humana prime sobre la generación estocástica de textos artificiales.

\newpage

\subsubsection{Evidencia Visual y Matriz Comparativa de Marcos de Competencia}

La articulación entre los marcos ciudadanos y pedagógicos descritos se sintetiza en la Figura 2.1, que modela la convergencia entre las cinco áreas fundamentales de DigComp 3.0, los dominios especializados de DigCompEdu y la capa transversal de Inteligencia Artificial que impregna la educación contemporánea.

\begin{figure}[htpb!]
\centering
\begin{tcolorbox}[colback=white, colframe=unmsm_gold, arc=3pt, boxrule=0.8pt, top=0.3cm, bottom=0.3cm, left=0.5cm, right=0.5cm]
\centering
\small\sffamily
\textbf{Figura 2.1.} Matriz Arquitectónica Integrada de Competencias Digitales e Inteligencia Artificial en Educación Superior\\[0.3cm]
\resizebox{\linewidth}{!}{%
\begin{tikzpicture}[>=stealth, font=\footnotesize\sffamily]
    % Bloque Superior IA Transversal (y = 3.2)
    \node [rectangle, draw=unmsm_red, fill=bg_arena, rounded corners=3pt, text width=12cm, text centered, inner sep=6pt] (ia_layer) at (0, 3.2) {
        \textbf{\textcolor{unmsm_red}{Capa Transversal de Inteligencia Artificial y Soberanía Epistémica (DigComp 3.0 / UNESCO 2024)}}\\
        Comprensión de LLMs $\cdot$ Auditoría de Alucinaciones $\cdot$ Prompting Crítico $\cdot$ Ética Algorítmica $\cdot$ Gobernanza de Datos
    };

    % Bloque DigComp 3.0 (Ciudadanía) (y = 0.4)
    \node [rectangle, draw=unmsm_red, fill=bg_marfil, rounded corners=3pt, text width=5.6cm, text centered, inner sep=5pt] (digcomp) at (-3.3, 0.4) {
        \textbf{\textcolor{unmsm_red}{DigComp 3.0: Ciudadanía Digital}}\\
        \scriptsize
        \begin{itemize}[leftmargin=1.1em, topsep=2pt, itemsep=1pt]
            \item 1. Búsqueda y gestión de datos
            \item 2. Comunicación y colaboración
            \item 3. Creación de contenidos digitales
            \item 4. Seguridad y protección integral
            \item 5. Resolución crítica de problemas
        \end{itemize}
    };

    % Bloque DigCompEdu (Docencia Universitaria) (y = 0.4)
    \node [rectangle, draw=unmsm_red, fill=bg_marfil, rounded corners=3pt, text width=5.6cm, text centered, inner sep=5pt] (digcompedu) at (3.3, 0.4) {
        \textbf{\textcolor{unmsm_red}{DigCompEdu: Pedagogía Universitaria}}\\
        \scriptsize
        \begin{itemize}[leftmargin=1.1em, topsep=2pt, itemsep=1pt]
            \item 1. Compromiso institucional
            \item 2. Curaduría y recursos digitales
            \item 3. Mediación del aprendizaje
            \item 4. Analíticas y evaluación formativa
            \item 5. Empoderamiento e inclusión (DUA)
            \item 6. Facilitación de la competencia
        \end{itemize}
    };

    % Bloque Base: Praxis Universitaria (y = -3.2)
    \node [rectangle, draw=unmsm_darkred, fill=bg_arena, rounded corners=3pt, text width=12cm, text centered, inner sep=5pt] (praxis) at (0, -3.2) {
        \textbf{\textcolor{unmsm_darkred}{Praxis Universitaria UNMSM: Agencia Crítica, Ética Científica y Equidad Social}}\\
        Trazabilidad DOI $\cdot$ Inoculación contra Desinformación $\cdot$ Accesibilidad Universal WCAG 2.2 $\cdot$ Criptografía NIST
    };

    % Conectores
    \draw [->, thick, draw=unmsm_red] (ia_layer) -- (digcomp);
    \draw [->, thick, draw=unmsm_red] (ia_layer) -- (digcompedu);
    \draw [->, thick, draw=unmsm_darkred] (digcomp) -- (praxis);
    \draw [->, thick, draw=unmsm_darkred] (digcompedu) -- (praxis);
\end{tikzpicture}%
}\\[0.2cm]
\scriptsize \textbf{Fuente:} Elaboración propia basada en Cosgrove y Cachia (2025), Vuorikari et al. (2022) y Redecker (2017).
\end{tcolorbox}
\end{figure}

Para clarificar el alcance operativo de cada marco de referencia y evitar solapamientos conceptuales, la Tabla 2.1 sistematiza una comparación formal entre DigComp 2.2, DigComp 3.0 y DigCompEdu, desglosando sus dimensiones, descriptores y mecanismos de verificación pedagógica.

\begin{table}[htpb!]
\centering
\small
\begin{tabularx}{\textwidth}{p{2.8cm} p{3.8cm} p{4.2cm} X}
\toprule
\textbf{Marco de Referencia} & \textbf{Audiencia y Propósito} & \textbf{Estructura y Dimensiones} & \textbf{Integración de Inteligencia Artificial} \\
\midrule
\textbf{DigComp 2.2}\newline(Vuorikari et al., 2022) & Ciudadanía general y fuerza laboral. Estandarización de habilidades ciudadanas en la UE. & 5 áreas transversales, 21 subcompetencias, 8 niveles de dominio y más de 250 descriptores prácticos. & Tratamiento preliminar: ejemplos descriptivos aislados sobre reconocimiento de asistentes de voz y sistemas automatizados simples. \\
\midrule
\textbf{DigComp 3.0}\newline(Cosgrove \& Cachia, 2025) & Ciudadanía en la era algorítmica. Respuestas a la automatización masiva y los LLMs. & 5 áreas clásicas reformuladas bajo mediación algorítmica y gobernanza de la atención digital. & Integración transversal obligatoria: prompts críticos, detección de alucinaciones, sesgos algorítmicos y auditoría de \textit{deepfakes}. \\
\midrule
\textbf{DigCompEdu}\newline(Redecker, 2017) & Docentes de todos los niveles educativos (primaria hasta posgrado universitario). & 6 áreas pedagógicas: profesional, recursos, enseñanza, evaluación, inclusión y competencia discente. & Foco en la orquestación didáctica: analíticas de aprendizaje formativas, tutores adaptativos y diseño centrado en lo humano. \\
\bottomrule
\end{tabularx}
\vspace{0.2cm}
\begin{flushleft}
\scriptsize \textbf{Nota.} Adaptado de Cosgrove y Cachia (2025), Vuorikari et al. (2022) y Redecker (2017).
\end{flushleft}
\end{table}

\newpage

\subsection{Curaduría Crítica de la Información, Inoculación Psicológica y Mitigación de la Desinformación}

\subsubsection{Arquitectura de la Duda Metódica: Preguntas Heurísticas y Trazabilidad Epistémica}

En la ecología informativa digital contemporánea, caracterizada por la infoxicación y la diseminación acelerada de contenidos generados automáticamente, la adquisición de certezas científicas exige la institucionalización de una \textbf{duda metódica procedimental}. La simple lectura pasiva de un titular o la consulta apresurada a un modelo generativo de lenguaje no constituye un acto de investigación académica.

La curaduría rigurosa de fuentes demanda la aplicación sistemática de preguntas heurísticas de control antes de incorporar cualquier dato a una producción académica o profesional:
\begin{itemize}
    \item \textbf{Criterio de Autoría y Filiación Institucional:} \textquestiondown Quién firma la afirmación? \textquestiondown Cuenta el autor con trayectoria verificable en el área de especialidad? \textquestiondown Cuál es su filiación institucional (universidad, instituto de investigación, organismo gubernamental)?
    \item \textbf{Criterio de Evidencia y Soporte Empírico:} \textquestiondown Qué metodología o conjunto de datos respalda la aseveración? \textquestiondown Se describen el tamaño de muestra, los métodos de recolección y las pruebas estadísticas empleadas?
    \item \textbf{Criterio de Temporalidad y Actualización:} \textquestiondown Cuándo fue publicado originalmente el estudio? \textquestiondown Ha sido modificado, corregido o refutado por investigaciones posteriores de mayor envergadura?
    \item \textbf{Criterio de Trazabilidad y Acceso a la Fuente Primaria:} \textquestiondown El documento proporciona enlaces directos a los registros de origen o se limita a citar resúmenes fragmentarios de terceros?
\end{itemize}

El componente medular de la trazabilidad científica moderna es el \textbf{Identificador de Objeto Digital} (\textit{Digital Object Identifier}, \textbf{DOI}). Administrado por la Fundación Internacional DOI (IDF) y agencias de registro como Crossref y DataCite, el DOI es un identificador alfanumérico persistente e interoperable que resuelve de forma permanente la ubicación de un documento en la web, inmune a los enlaces rotos originados por la migración de servidores o el cambio de plataformas editoriales.

La anatomía formal de un DOI consta invariablemente de dos segmentos separados por una barra inclinada:
\begin{center}
\vspace{0.15cm}
\begin{tikzpicture}[font=\small\sffamily]
    \node [rectangle, draw=unmsm_red, fill=bg_marfil, rounded corners=3pt, inner sep=6pt] (pref) at (-3.2, 0) {\textbf{\texttt{10.1038}}};
    \node at (0, 0) {\Large\textbf{/}};
    \node [rectangle, draw=unmsm_darkred, fill=bg_arena, rounded corners=3pt, inner sep=6pt] (suf) at (3.5, 0) {\textbf{\texttt{s41467-022-30073-5}}};
    \node [font=\scriptsize\sffamily, text=unmsm_red, below=0.15cm of pref] {\textbf{Prefijo}: Agencia y Editorial (Crossref)};
    \node [font=\scriptsize\sffamily, text=unmsm_darkred, below=0.15cm of suf] {\textbf{Sufijo}: Objeto Específico e Identificador de Manuscrito};
\end{tikzpicture}
\vspace{0.15cm}
\end{center}
La resolución del DOI a través de su protocolo estándar (\url{https://doi.org/10.1038/s41467-022-30073-5}) garantiza el acceso instantáneo a los metadatos oficiales del artículo en el servidor de la editorial, lo que permite a los investigadores contrastar la autenticidad del manuscrito y protegerse frente al fenómeno del secuestro de revistas (\textit{journal hijacking}) y la proliferación de sitios clonados que suplantan publicaciones prestigiosas para capturar pagos ilegítimos.

\subsubsection{Psicología Cognitiva de la Desinformación: Pausas de Precisión (Accuracy Prompts)}

Uno de los enigmas más desconcertantes de la era digital es por qué millones de personas instruidas comparten cotidianamente noticias falsas o información médica errónea en sus redes sociales, incluso cuando poseen el conocimiento necesario para advertir la inconsistencia si se les interroga de manera directa.

En una investigación experimental decisiva publicada en \textit{Nature Communications}, Gordon Pennycook y David Rand (2022) revelaron el mecanismo cognitivo subyacente. Los autores demuestran que la arquitectura sociotécnica de las redes sociales está deliberadamente optimizada para recompensar el enganche afectivo, la inmediatez y la validación social (reacciones, retuits, comentarios). En ese entorno hiperestimulado, la atención del usuario se focaliza de forma casi exclusiva en el valor social y la resonancia moral del mensaje, desactivando de facto la deliberación epistémica. En consecuencia, las personas comparten información falsa no por una incapacidad intelectual intrínseca ni por un deseo consciente de mentir, sino por una simple \textbf{distracción atencional}: el entorno no les solicita pensar en la exactitud.

Como solución a esta patología atencional, Pennycook y Rand (2022) formularon la intervención de las \textbf{Pausas de Precisión} (\textit{Accuracy Prompts}). Consiste en un estímulo contextual mínimo y replicable ---por ejemplo, solicitar al usuario calificar la exactitud de un titular neutro antes de iniciar su sesión de navegación--- que redirecciona temporalmente el foco de la memoria de trabajo hacia la dimensión de veracidad. 

La evidencia cuantitativa obtenida por los autores a través de múltiples ensayos controlados aleatorizados confirma que esta simple intervención sutil (\textit{nudge}) reduce entre un 30\% y un 50\% la predisposición a compartir información falsa en plataformas digitales, duplicando la tasa de discernimiento epistémico. Para la educación universitaria, este hallazgo implica que la enseñanza de la alfabetización digital debe incorporar rutinas reflexivas explícitas: inducir intencionalmente un instante de pausa y formulación de preguntas de comprobación antes de que el estudiante o docente redistribuya un enlace o adopte un resultado algorítmico en su práctica académica.

\subsubsection{Teoría de la Inoculación Psicológica y Pre-bunking frente a Técnicas de Manipulación}

Aunque la verificación posterior de hechos (\textit{fact-checking} o desmentido reactivo) es necesaria, presenta limitaciones estructurales severas: la desinformación se propaga a una velocidad exponencialmente superior al tiempo que insume la verificación humana, y una vez que una falsedad es internalizada en los esquemas cognitivos, los desmentidos suelen generar el efecto de rebote (\textit{backfire effect}) o dejar secuelas de desinformación residual (\textit{continued influence effect}).

Para superar estas limitaciones, Jon Roozenbeek, Sander van der Linden, Beth Goldberg, Steve Rathje y Stephan Lewandowsky (2022), en un estudio histórico publicado en \textit{Science Advances}, desarrollaron y validaron a escala masiva la teoría de la \textbf{Inoculación Psicológica} o \textbf{Pre-bunking}. Inspirada en la metáfora biológica de las vacunas médicas, la inoculación psicológica postula que exponer de manera controlada y preventiva a las personas a microdosis de las técnicas retóricas y argumentativas utilizadas por los agentes desinformadores genera anticuerpos cognitivos que aumentan la resistencia psicológica ante ataques futuros.

Roozenbeek et al. (2022) identificaron cinco técnicas capitales de manipulación sistemática:
\begin{enumerate}
    \item \textbf{Apelación a Emociones Hiperbólicas:} Uso deliberado de palabras con alta carga afectiva (indignación moral, miedo, euforia) para suspender el análisis crítico del lector.
    \item \textbf{Polarización Dicotómica (\textit{In-group vs. Out-group}):} Construcción artificial de falsos dilemas que enfrentan de forma irreconciliable a colectivos identitarios, culpabilizando a un chivo expiatorio de problemas estructurales complejos.
    \item \textbf{Falsos Expertos:} Presentación de testimonios o credenciales académicas en áreas no pertinentes para respaldar afirmaciones pseudocientíficas carentes de consenso en la comunidad de especialistas.
    \item \textbf{Narrativas Conspirativas:} Proposición de explicaciones monistas e infalsables que atribuyen eventos fortuitos a confabulaciones secretas de élites todopoderosas, descartando cualquier evidencia contraria como parte del complot.
    \item \textbf{Descalificación \textit{Ad Hominem}:} Ataque sistemático a la integridad moral o rasgos personales del interlocutor para desviar la atención de la validez o invalidez de sus argumentos lógicos.
\end{enumerate}

La Tabla 2.2 sintetiza estas técnicas de manipulación, sus mecanismos cognitivos de penetración y las estrategias operativas de contra-manipulación que deben enseñarse en las aulas universitarias.

\begin{table}[htpb!]
\centering
\small
\begin{tabularx}{\textwidth}{p{3.0cm} p{3.2cm} p{4.2cm} X}
\toprule
\textbf{Técnica de Manipulación} & \textbf{Mecanismo Psicológico} & \textbf{Manifestación Típica en Medios} & \textbf{Estrategia de Inoculación (Pre-bunking)} \\
\midrule
\textbf{Apelación Emocional} & Sesgo de saliencia afectiva y secuestro de la atención. & Titulares alarmistas con mayúsculas y adjetivos extremos (\textit{"¡Descubrimiento aterrador oculta la verdad!"}). & Identificar el léxico hiperbólico y aplicar una pausa reflexiva de enfriamiento emocional. \\
\midrule
\textbf{Falsos Expertos} & Efecto halo y apelación falaz a la autoridad (\textit{ad verecundiam}). & Citar a graduados de áreas inconexas para opinar sobre epidemiología o criptografía avanzada. & Verificar el factor de impacto del autor en Scopus/WoS y su historial en revistas con arbitraje ciego. \\
\midrule
\textbf{Falsa Dicotomía} & Pensamiento maniqueo y polarización de grupo. & Presentación forzada de opciones mutuamente excluyentes (\textit{"O estás con el progreso o estás contra él"}). & Desmontar el dilema identificando matices intermedios y terceras vías fundamentadas empíricamente. \\
\midrule
\textbf{Narrativa Conspirativa} & Ilusión de causalidad y necesidad de cierre cognitivo. & Hipótesis omniexplicativas que tachan toda evidencia científica en contra como "complicidad con el sistema". & Contrastar con el principio de la Navaja de Ockham y exigir pruebas extraordinarias a afirmaciones extraordinarias. \\
\midrule
\textbf{Ataque \textit{Ad Hominem}} & Desvío de atención del argumento a la persona. & Descalificar un informe metodológico insultando la apariencia o ideología política de sus autores. & Reenfocar el debate estrictamente en la consistencia de los datos y la reproducibilidad metodológica. \\
\bottomrule
\end{tabularx}
\vspace{0.2cm}
\begin{flushleft}
\scriptsize \textbf{Nota.} Adaptado de Roozenbeek et al. (2022) y Valverde-Berrocoso et al. (2022).
\end{flushleft}
\end{table}

\newpage

\subsection{Creación Inclusiva de Contenidos, Accesibilidad Universal (WCAG 2.2) y Seguridad de la Identidad Digital}

\subsubsection{Ingeniería de la Accesibilidad Web y Documental: Principios POUR y Evidencia Empírica}

En la universidad moderna, la creación y publicación de contenidos digitales no puede constituir una actividad discrecional sujeta al arbitrio estético del docente o investigador. La educación superior como bien público exige el cumplimiento irrestricto del derecho a la accesibilidad universal para toda la comunidad académica, incluyendo a personas con discapacidades visuales, auditivas, motoras o neurodivergencias.

El estándar técnico global de referencia son las Pautas de Accesibilidad para el Contenido Web (\textit{Web Content Accessibility Guidelines}, versión \textbf{WCAG 2.2}), formuladas por el Consorcio World Wide Web (W3C). Como documentan rigurosamente Ara, Sik-Lanyi y Kelemen (2024) en su revisión sistemática sobre ingeniería de accesibilidad en los procesos de evaluación web, la conformidad con WCAG 2.2 se articula en torno a cuatro principios fundacionales conocidos por el acrónimo \textbf{POUR}:
\begin{enumerate}
    \item \textbf{Perceptible (Perceivable):} La información y los componentes de la interfaz deben presentarse a los usuarios de manera que puedan ser captados por sus sentidos disponibles. Exige que todo elemento no textual (imágenes, gráficos vectoriales, diagramas) cuente con una alternativa textual concisa y descriptiva (\texttt{alt text}), que el contraste cromático entre texto y fondo cumpla un umbral matemático mínimo de $4.5:1$ para texto estándar y $3:1$ para títulos grandes, y que los videos incorporen subtitulación sincronizada y transcripciones completas.
    \item \textbf{Operable (Operable):} La totalidad de las funciones de navegación e interacción debe ser operable mediante diferentes dispositivos periféricos, especialmente el teclado estándar, sin exigir movimientos de motricidad fina ni atrapar al usuario en "trampas de foco" (\textit{keyboard traps}). Los tiempos de interacción deben ser configurables y los elementos clicables deben tener un área mínima de toque de $24 \times 24$ píxeles CSS para prevenir activaciones accidentales.
    \item \textbf{Comprensible (Understandable):} El contenido textual debe ser legible, estructurado con jerarquías lógicas (\texttt{H1}, \texttt{H2}, \texttt{H3}) y redactado siguiendo directrices de lectura fácil. Los formularios y las plataformas de trámites académicos deben proveer instrucciones inequívocas, mecanismos predictivos de prevención de errores y mensajes claros para corregir campos defectuosos sin borrar la información ingresada.
    \item \textbf{Robusto (Robust):} El código y los documentos deben ser interpretables de forma fiable por una amplia diversidad de agentes de usuario, asegurando total interoperabilidad con tecnologías asistivas (como lectores de pantalla NVDA, JAWS o VoiceOver, magnificadores de caracteres y conmutadores de hardware).
\end{enumerate}

Ara et al. (2024) subrayan que la accesibilidad no debe ser concebida como un añadido cosmético de última hora, sino como una disciplina de ingeniería integrada desde la concepción del diseño curricular y arquitectónico. En entornos universitarios, esto demanda que las aulas virtuales (Moodle, Canvas), los repositorios institucionales y las guías de estudio en PDF sean auditados sistemáticamente contra los criterios de conformidad nivel AA de las WCAG 2.2.

\subsubsection{Seguridad de la Identidad Digital y Autenticación Criptográfica: NIST SP 800-63B-4}

La digitalización integral de la gestión universitaria (expedientes de calificaciones, registros de matrícula, propiedad intelectual de tesis y pasarelas de pago) convierte a las credenciales académicas en blancos de alto valor para cibercriminales y atacantes informáticos. Durante años, las instituciones educativas impusieron políticas tradicionales obsoletas basadas en exigir a los usuarios contraseñas complejas y forzar su rotación arbitraria cada noventa días.

En la reciente actualización de sus directrices federales, el Instituto Nacional de Estándares y Tecnología de los Estados Unidos (NIST), a través del estándar \textbf{NIST SP 800-63B-4} redactado por David Temoshok y colaboradores (2025), demolió de forma concluyente estos paradigmas heredados. Temoshok et al. (2025) demuestran que las políticas de rotación periódica forzada de contraseñas son contraproducentes para la seguridad: los usuarios, abrumados por la fatiga nemotécnica, tienden a aplicar modificaciones mínimas y predecibles (como cambiar un número correlativo final o añadir un signo de exclamación), generando patrones sistemáticos fácilmente explotables por ataques de fuerza bruta y diccionarios basados en redes neuronales.

El estándar NIST SP 800-63B-4 redefine la autenticación digital segura a través de los siguientes lineamientos:
\begin{itemize}
    \item \textbf{Superación de Contraseñas Estáticas Simples:} Promoción activa de contraseñas largas (\textit{passphrases}) de más de quince caracteres basadas en oraciones memorables, eliminando exigencias arbitrarias de caracteres especiales que aumentan la frustración del usuario.
    \item \textbf{Verificación contra Listas de Credenciales Comprometidas:} Obligación del servidor de cotejar en tiempo real las contraseñas elegidas por los usuarios contra bases de datos mundiales de credenciales filtradas en violaciones de datos (\textit{breaches}).
    \item \textbf{Autenticación Multifactorial Resistente al Phishing (FIDO2 / WebAuthn):} Demostración de que el segundo factor basado en mensajes SMS o códigos OTP generados en aplicaciones móviles es vulnerable a ataques de intermediario (\textit{Adversary-in-the-Middle}, AitM), intercambio fraudulento de SIM (\textit{SIM swapping}) e ingeniería social. La norma NIST exige la transición hacia factores criptográficos basados en el estándar FIDO2 / WebAuthn (\textbf{passkeys}), donde la autenticación se realiza mediante pares de claves asimétricas enlazadas criptográficamente al dominio del servidor, haciendo matemáticamente inviable el robo de credenciales mediante phishing.
\end{itemize}

\subsubsection{Patrones Comerciales Oscuros (Dark Patterns) y Datificación en Plataformas Educativas}

En paralelo a las amenazas de seguridad informática, la comunidad académica enfrenta el riesgo invisible de los \textbf{patrones comerciales oscuros} (\textit{dark commercial patterns}) documentados exhaustivamente por la OCDE / OECD (2022). Los patrones oscuros son diseños de interfaz intencionalmente manipulativos que explotan sesgos cognitivos y heurísticas humanas para inducir a los usuarios a tomar decisiones que favorecen los intereses comerciales del proveedor de la plataforma en detrimento de su propia privacidad o bienestar financiero.

Entre las tipologías más prevalentes en plataformas educativas comerciales y herramientas freemium adoptadas informalmente en universidades destacan:
\begin{enumerate}
    \item \textbf{Consentimiento Forzado u Obscurecido:} Interfaces donde el botón para aceptar el rastreo masivo de cookies de terceros y telemetría de comportamiento es grande, colorido y visible, mientras que la opción para rechazarlo o configurarlo se oculta tras múltiples menús en gris claro sobre blanco.
    \item \textbf{Confirmación Culposa (\textit{Confirm-shaming}):} Uso de textos manipulativos que inducen culpabilidad o vergüenza emocional al rechazar una opción de pago o actualización (ej. \textit{"No gracias, prefiero no mejorar mi rendimiento académico"}).
    \item \textbf{Laberintos de Cancelación (\textit{Roach Motel}):} Procedimientos donde suscribirse a un servicio requiere un solo clic, pero darse de baja o eliminar la cuenta exige completar llamadas telefónicas, enviar formularios físicos o atravesar laberintos burocráticos.
\end{enumerate}

La Tabla 2.3 sintetiza las dimensiones de accesibilidad WCAG 2.2, los controles criptográficos NIST SP 800-63B-4 y las directrices de mitigación de patrones oscuros para entornos educativos.

\begin{table}[htpb!]
\centering
\small
\begin{tabularx}{\textwidth}{p{3.2cm} p{3.8cm} p{4.0cm} X}
\toprule
\textbf{Área Normativa} & \textbf{Estándar de Referencia} & \textbf{Riesgo / Vulnerabilidad} & \textbf{Control Técnico e Institucional} \\
\midrule
\textbf{Accesibilidad Web y Documental} & WCAG 2.2 (Nivel AA)\newline(Ara et al., 2024) & Exclusión de estudiantes con discapacidad visual o motriz; rechazo de lectores de pantalla. & Implementación de principios POUR: contraste $\ge 4.5:1$, navegación por teclado, texto alternativo riguroso. \\
\midrule
\textbf{Autenticación e Identidad} & NIST SP 800-63B-4\newline(Temoshok et al., 2025) & Secuestro de cuentas por phishing AitM y reciclaje de contraseñas filtradas. & Despliegue de claves criptográficas FIDO2/WebAuthn (\textit{passkeys}) y contraseñas de alta entropía sin caducidad forzada. \\
\midrule
\textbf{Privacidad y Datificación} & OCDE / OECD (2022) / RGPD / Ley 29733 (Perú) & Extracción no consentida de datos biométricos y hábitos de estudio en plataformas EdTech. & Prohibición de patrones oscuros de consentimiento; almacenamiento seguro en servidores institucionales soberanos. \\
\bottomrule
\end{tabularx}
\vspace{0.2cm}
\begin{flushleft}
\scriptsize \textbf{Nota.} Adaptado de Ara et al. (2024), Temoshok et al. (2025) y OCDE / OECD (2022).
\end{flushleft}
\end{table}

\newpage

\subsection{Alfabetización Crítica en Inteligencia Artificial (AI Literacy) y Sistemas de Aprendizaje Adaptativo}

\subsubsection{El Modelo de Cuatro Dimensiones de la Alfabetización en IA}

La vertiginosa irrupción de la inteligencia artificial generativa ha transformado la relación entre los aprendices y el conocimiento. Frente a visiones tecnocéntricas que conciben la alfabetización en IA como el mero adiestramiento en la ingeniería de prompts o en la sintaxis de lenguajes de programación, Davy Tsz Kit Ng, Jac Ka Lok Leung, Samuel Kai Wah Chu y Maggie Shen Qiao (2021) formularon un influyente marco teórico multidimensional que delimita los componentes indispensables de la \textbf{Alfabetización en Inteligencia Artificial} (\textit{AI Literacy}).

A partir de una rigurosa revisión exploratoria, Ng et al. (2021) estructuran el constructo en cuatro dimensiones conceptuales interconectadas:
\begin{enumerate}
    \item \textbf{Conocer y Comprender la IA (\textit{Know \& Understand}):} Comprensión básica de qué es la inteligencia artificial, distinguiendo entre sistemas basados en reglas lógicas deterministas y algoritmos de aprendizaje automático estadístico (\textit{Machine Learning}). Abarca la comprensión de cómo los modelos procesan datos multimodales (texto, audio, visión computacional), qué representa el entrenamiento supervisado y no supervisado, y cuáles son las limitaciones físicas y algorítmicas de los modelos predictivos.
    \item \textbf{Usar y Aplicar la IA (\textit{Use \& Apply}):} Habilidad técnica para seleccionar, interactuar y orquestar herramientas de inteligencia artificial en la resolución de problemas académicos y profesionales específicos. Incluye la formulación de directivas estructuradas (\textit{prompt engineering}), la contextualización semántica de consultas y la integración de salidas de IA en canalizaciones de trabajo complejas.
    \item \textbf{Evaluar y Crear con IA (\textit{Evaluate \& Create}):} Capacidad de orden superior para auditar críticamente los resultados generados por los modelos. Involucra la detección de alucinaciones matemáticas y conceptuales, la identificación de sesgos de representación, la verificación de coherencia argumental frente a fuentes primarias y la capacidad de co-crear soluciones innovadoras combinando la heurística humana con la velocidad de cómputo algorítmica.
    \item \textbf{Ética e Impacto de la IA (\textit{AI Ethics}):} Evaluación reflexiva de las implicancias éticas, axiológicas y jurídicas de los sistemas automatizados. Analiza la propiedad intelectual de las creaciones sintéticas, la privacidad de los datos suministrados a los modelos, la equidad algorítmica, el desplazamiento laboral cualificado y la huella de carbono de los centros de procesamiento de inferencia masiva.
\end{enumerate}

Un aporte epistemológico fundamental de Ng et al. (2021) es la urgencia de desmitificar el \textbf{modelo mental antropomórfico}. Los modelos de lenguaje masivo contemporáneos no "comprenden" el significado intrínseco de sus respuestas ni poseen conciencia ni intencionalidad; operan como descomunales sistemas de cálculo estadístico que predicen la distribución de probabilidad del siguiente token textual en función de patrones inferidos de su corpus de entrenamiento. Interpretar sus respuestas bajo categorías de psicología humana constituye un error categorial que induce a una peligrosa e infundada sobreconfianza epistémica.

\subsubsection{Validación Psicométrica de Competencias en IA: La Escala MAILS}

Para operacionalizar empíricamente la medición de estas competencias y permitir el diagnóstico riguroso en poblaciones universitarias, Astrid Carolus, Maria J. Koch, Sarah Straka, Marc Erich Latoschik y Carolin Wienrich (2023) desarrollaron y validaron psicométricamente la \textbf{Escala Meta de Alfabetización en IA} (\textit{Meta AI Literacy Scale}, \textbf{MAILS}), publicada en \textit{Computers in Human Behavior: Artificial Humans}.

Carolus et al. (2023) identificaron una paradoja psicológica recurrente en la adopción tecnológica: la brecha entre la competencia autopercibida (\textit{perceived competence}) y el rendimiento operacional comprobable (\textit{actual performance}). Con frecuencia, los usuarios que reportan niveles elevados de familiaridad con herramientas de IA incurren en graves fallos de juicio debido a la complacencia cognitiva generada por la elocuencia superficial de los modelos generativos (el llamado efecto de adulación o \textit{sycophancy}).

La escala MAILS supera esta deficiencia articulando cuatro subescalas de alta confiabilidad psicométrica (coeficientes $\alpha$ de Cronbach $\ge 0.82$):
\begin{itemize}
    \item \textbf{Subescala 1: Comprensión Funcional de la IA:} Evalúa el conocimiento sobre los mecanismos operativos de los algoritmos y la capacidad para distinguir entre capacidades reales y limitaciones técnicas.
    \item \textbf{Subescala 2: Habilidad de Interacción y Uso:} Mide la destreza procedimental para modular los prompts, iterar conversacionalmente y guiar a la IA hacia los objetivos deseados.
    \item \textbf{Subescala 3: Autorregulación y Evaluación Crítica Metacognitiva:} Monitorea la predisposición del usuario a cuestionar sistemáticamente las salidas del modelo, comprobar datos de forma independiente y detectar alucinaciones lógicas.
    \item \textbf{Subescala 4: Conciencia Ética y Social:} Cuantifica la sensibilidad del individuo ante dilemas de privacidad, sesgos discriminatorios y dilemas de autoría académica.
\end{itemize}

\subsubsection{Bienestar Digital, Carga Cognitiva y Desconexión: Más Allá del Mito del Digital Detox}

El incremento de la mediación tecnológica en la vida académica ha generado preocupación social sobre sus impactos en la salud mental, el descanso y la concentración profunda. En el debate público se ha popularizado la noción del "detox digital" (\textit{digital detox}) como una panacea consistente en la desconexión total temporal de todo dispositivo.

No obstante, la investigación científica arroja un panorama sustancialmente más matizado. En una rigurosa revisión sistemática de la literatura publicada en \textit{Mobile Media \& Communication}, Theda Radtke, Tamara Apel, Katja Schenkel, Jan Keller y Eike von Lindern (2022) analizaron los efectos empíricamente documentados de las intervenciones de desconexión digital. Sus hallazgos demuestran que las prescripciones de desconexión total obligatoria producen resultados contradictorios e inestables: si bien pueden reducir temporalmente las distracciones superficiales, frecuentemente generan ansiedad inducida por el aislamiento social (\textit{Fear of Missing Out}, FOMO) y provocan un efecto rebote de hiperconsumo al concluir el período de abstinencia, sin solucionar las causas estructurales del estrés laboral o académico.

Radtke et al. (2022) postulan que la alfabetización digital avanzada debe promover un enfoque de \textbf{bienestar digital sostenible} basado en el rediseño consciente de las condiciones de interacción, más que en la abstinencia punitiva. Esto implica:
\begin{enumerate}
    \item \textbf{Intencionalidad de Uso frente a la Reactividad Pasiva:} Pasar de un consumo compulsivo conducido por algoritmos de notificación a sesiones deliberadas de trabajo con propósitos definidos.
    \item \textbf{Arquitectura del Descanso y Ergonomía Temporal:} Delimitación estricta de franjas horarias libres de pantallas, desactivación de notificaciones push de aplicaciones no críticas y preservación de la higiene del sueño mediante la eliminación de dispositivos emisores de luz azul antes de dormir.
    \item \textbf{Gestión de la Carga Cognitiva y Trabajo Profundo (\textit{Deep Work}):} Bloques estructurados de estudio sin interrupciones digitales para permitir la consolidación en la memoria a largo plazo, complementados por una alfabetización socioemocional en redes (Polanco-Levicán \& Salvo-Garrido, 2022).
\end{enumerate}

\newpage

\subsubsection{Esquema Vectorial y Matriz Diagnóstica MAILS}

La articulación entre el modelo de cuatro dimensiones de Ng et al. (2021), los constructos psicométricos de MAILS (Carolus et al., 2023) y la supervisión humana se modela en la Figura 2.2.

\begin{figure}[htpb!]
\centering
\begin{tcolorbox}[colback=white, colframe=unmsm_gold, arc=3pt, boxrule=0.8pt, top=0.3cm, bottom=0.3cm, left=0.5cm, right=0.5cm]
\centering
\small\sffamily
\textbf{Figura 2.2.} Arquitectura de Alfabetización Crítica en IA y Ciclo Metacognitivo MAILS\\[0.3cm]
\begin{tikzpicture}[node distance=1.4cm, auto, >=stealth, font=\footnotesize\sffamily]
    % Nodo Central: Sujeto Crítico
    \node [rectangle, draw=unmsm_red, fill=bg_arena, rounded corners=3pt, text width=4.8cm, text centered, inner sep=6pt] (sujeto) {
        \textbf{\textcolor{unmsm_red}{Sujeto Académico}}\\[0.1cm]
        \scriptsize Agencia Epistémica $\cdot$ Metacognición
    };

    % Nodo Superior: Comprensión y Uso
    \node [rectangle, draw=unmsm_red, fill=bg_marfil, rounded corners=3pt, text width=4.5cm, text centered, above=1.2cm of sujeto, inner sep=5pt] (know_use) {
        \textbf{\textcolor{unmsm_red}{1. Comprender y Usar (Ng)}}\\
        \scriptsize Modelos probabilísticos $\cdot$ Prompting
    };

    % Nodo Derecho: Sistema de IA
    \node [rectangle, draw=unmsm_darkred, fill=bg_marfil, rounded corners=3pt, text width=4.5cm, text centered, right=1.6cm of sujeto, inner sep=5pt] (modelo_ia) {
        \textbf{\textcolor{unmsm_darkred}{Sistema Algorítmico (LLM)}}\\
        \scriptsize Predicción estocástica $\cdot$ Alucinaciones
    };

    % Nodo Inferior: Evaluación y Ética
    \node [rectangle, draw=unmsm_red, fill=bg_marfil, rounded corners=3pt, text width=4.5cm, text centered, below=1.2cm of sujeto, inner sep=5pt] (eval_ethics) {
        \textbf{\textcolor{unmsm_red}{2. Evaluar y Ética (MAILS)}}\\
        \scriptsize Detección de sesgos $\cdot$ Trazabilidad
    };

    % Flechas del Ciclo
    \draw [->, thick, draw=unmsm_red] (sujeto) -- node [above] {\scriptsize Directiva} (know_use);
    \draw [->, thick, draw=unmsm_red] (know_use) -| node [above, pos=0.25] {\scriptsize Prompt contextual} (modelo_ia);
    \draw [->, thick, draw=unmsm_darkred] (modelo_ia) |- node [below, pos=0.75] {\scriptsize Salida sintética} (eval_ethics);
    \draw [->, thick, draw=unmsm_red] (eval_ethics) -- node [below] {\scriptsize Supervisión crítica} (sujeto);
\end{tikzpicture}\\[0.2cm]
\scriptsize \textbf{Fuente:} Elaboración propia basada en Ng et al. (2021) y Carolus et al. (2023).
\end{tcolorbox}
\end{figure}

Para permitir su aplicación diagnóstica inmediata en las facultades universitarias, la Tabla 2.4 desglosa los reactivos muestra de la escala MAILS y sus correspondientes estrategias pedagógicas de remediación.

\begin{table}[htpb!]
\centering
\small
\begin{tabularx}{\textwidth}{p{3.2cm} p{4.2cm} p{3.8cm} X}
\toprule
\textbf{Subescala MAILS} & \textbf{Reactivo Muestra (Ítem Likert 1--6)} & \textbf{Constructo Psicológico} & \textbf{Intervención Pedagógica Remediadora} \\
\midrule
\textbf{1. Comprensión Funcional} & \textit{"Comprendo que las respuestas de un modelo de lenguaje son predicciones probabilísticas y no hechos verificados."} & Modelo mental no antropomórfico de la computación. & Laboratorio práctico de deconstrucción algorítmica: generación de continuaciones probabilísticas disímiles con distintas temperaturas. \\
\midrule
\textbf{2. Habilidad de Uso e Interacción} & \textit{"Soy capaz de reformular iterativamente una directiva cuando la herramienta de IA entrega un resultado ambiguo o incompleto."} & Autoeficacia procedimental en ingeniería de prompts. & Talleres de prompting estructurado con marcos de rol, restricciones negativas y adición de contexto de dominio. \\
\midrule
\textbf{3. Evaluación y Metacognición} & \textit{"Cotejo de forma rutinaria las afirmaciones y citas generadas por IA contra fuentes primarias indexadas con DOI."} & Vigilancia epistémica y mitigación de sobreconfianza. & Auditoría ciega de textos sintéticos: detección guiada de citas inventadas y alucinaciones en papers reales. \\
\midrule
\textbf{4. Conciencia Ética y Social} & \textit{"Considero las implicancias de privacidad y cesión de derechos antes de subir documentos inéditos a plataformas externas."} & Conciencia bioética, jurídica y de gobernanza de datos. & Análisis de contratos de términos de servicio y diseño de políticas de gobernanza de datos en laboratorios. \\
\bottomrule
\end{tabularx}
\vspace{0.2cm}
\begin{flushleft}
\scriptsize \textbf{Nota.} Adaptado de Carolus et al. (2023) y Ng et al. (2021).
\end{flushleft}
\end{table}

\newpage

\subsection{Casos de Estudio Aplicados en Entornos Educativos y Comunitarios}

\subsubsection{Caso 1: Gestión de Identidad y Verificación Anti-Phishing en Trámites Estudiantiles (El Caso Rosa Adaptado)}

\subsubsection*{Contexto:}
Rosa es una estudiante universitaria de la Universidad Nacional Mayor de San Marcos que realiza actividades académicas y comunitarias. Durante el período de matrícula extemporánea y talleres extracurriculares, Rosa recibe un mensaje de mensajería instantánea reenviado a un grupo estudiantil no oficial. El mensaje anuncia una supuesta convocatoria de última hora para un "Taller de Reparación y Mantenimiento Tecnológico", prometiendo una certificación oficial universitaria con valor curricular y créditos extracurriculares.

El mensaje contiene un enlace acortado (\texttt{https://bit.ly/taller-reserva-2026}) y enfatiza un sentido de extrema urgencia: advierte que solo restan tres cupos disponibles y que, para asegurar la vacante, la persona interesada debe ingresar de inmediato al enlace, subir una fotografía por ambas caras de su Documento Nacional de Identidad (DNI), ingresar su usuario y contraseña del correo institucional \texttt{@unmsm.edu.pe}, y realizar un depósito bancario de 50 soles mediante una billetera digital a un número telefónico personal en concepto de "reserva de materiales".

\subsubsection*{Análisis e Implementación:}
En lugar de ceder a la urgencia emocional inducida por el mensaje (una manifestación arquetípica de manipulación afectiva para desactivar la deliberación analítica), Rosa activa el protocolo de curaduría crítica y seguridad digital estructurado en cuatro etapas:

\begin{enumerate}
    \item \textbf{Inspección Técnica de la URL y del Canal de Envío:} Rosa examina el enlace acortado en un entorno seguro antes de interactuar con él. Al resolver la redirección en un expansor de enlaces, constata que la URL de destino no pertenece al dominio oficial \texttt{unmsm.edu.pe}, sino a un servidor gratuito alojado en un dominio comercial sospechoso (\url{formulario-acreditaciones-peru.net}).
    \item \textbf{Contraste en Canales Oficiales Primarios:} Rosa accede directamente al portal web oficial de la Facultad de Ciencias Matemáticas de la UNMSM y al Sistema Único de Matrícula (SUM). En los canales formales constata que el taller existe como iniciativa de proyección social comunitaria, pero con condiciones diametralmente distintas: la inscripción es completamente gratuita, el único requisito es registrar el código de estudiante y correo institucional, y en ninguna circunstancia la universidad solicita contraseñas ni depósitos a cuentas individuales.
    \item \textbf{Contención y Alerta Comunitaria:} Rosa descarta de inmediato el enlace fraudulento. Acto seguido, redacta una advertencia respetuosa y precisa para el grupo de mensajería estudiantil, explicando con capturas de pantalla la discrepancia entre el enlace malicioso y el canal oficial de la universidad, evitando que sus compañeros entreguen datos personales o dinero. Asimismo, remite el incidente a la Oficina General de Tecnologías de Información (OGTI) de la UNMSM para el bloqueo perimetral del dominio falso.
    \item \textbf{Aseguramiento Criptográfico de la Identidad Digital:} Para blindar su propia cuenta ante futuros ataques de ingeniería social, Rosa ingresa a la configuración de seguridad de su cuenta institucional e implementa las directrices de NIST SP 800-63B-4: actualiza su contraseña por una frase de paso de 18 caracteres de alta entropía y activa la autenticación de doble factor mediante una clave criptográfica FIDO2 registrada en su dispositivo móvil.
\end{enumerate}

\subsubsection*{Interpretación y Resultados:}
El caso de Rosa evidencia de forma tangible la integración sinérgica entre la competencia técnica instrumental y el discernimiento crítico. La estudiante no necesitó ser especialista en ciberseguridad forense para neutralizar una amenaza de suplantación de identidad; bastó con la aplicación rigurosa de una pausa de precisión, la formulación de preguntas sobre el canal primario y el conocimiento de la arquitectura de dominios web. La gestión evitó la pérdida patrimonial y la filtración de credenciales sensibles, demostrando que la alfabetización digital es, en última instancia, una práctica de autocuidado y solidaridad comunitaria.

\newpage

\subsubsection{Caso 2: Estudio Empírico Cuantitativo de Alfabetización en IA en la UNMSM con Escala MAILS}

\subsubsection*{Contexto:}
Para evaluar empíricamente el impacto de un programa formativo formal en la adquisición de competencias críticas en inteligencia artificial, el Departamento Académico de la Facultad de Ciencias Matemáticas de la Universidad Nacional Mayor de San Marcos condujo un estudio experimental controlado durante el semestre académico 2026-I.

La cohorte de estudio estuvo conformada por una muestra probabilística estratificada de $N = 240$ participantes (200 estudiantes de pregrado y 40 docentes universitarios). Los participantes fueron asignados de manera aleatoria a dos ramas de estudio balanceadas:
\begin{itemize}
    \item \textbf{Grupo Experimental ($n = 120$):} Participó en el programa intensivo \textit{"Alfabetización Crítica en IA y Soberanía Epistémica"} de 8 semanas de duración (32 horas pedagógicas presenciales y virtuales), estructurado formalmente en los marcos DigCompEdu, Ng et al. (2021) y el protocolo de auditoría MAILS.
    \item \textbf{Grupo de Control ($n = 120$):} Continuó con las actividades académicas habituales sin intervención pedagógica específica en IA, utilizando herramientas tecnológicas según su criterio personal previo.
\end{itemize}

\subsubsection*{Análisis e Implementación:}
La evaluación cuasiexperimental contempló la administración de dos baterías diagnósticas en fases pre-test (semana 1) y post-test (semana 8):
\begin{enumerate}
    \item \textbf{Batería Psicométrica MAILS:} Cuestionario validado de 16 ítems en escala Likert (1 = Totalmente en desacuerdo, a 6 = Totalmente de acuerdo) midiendo las cuatro subescalas: Comprensión Funcional ($\text{CF}$), Habilidad de Uso ($\text{HU}$), Evaluación Crítica ($\text{EC}$) y Conciencia Ética ($\text{C\text{\'E}}$).
    \item \textbf{Prueba de Discernimiento Epistémico y Detección de Alucinaciones:} Prueba de rendimiento ciego consistente en la lectura de 20 textos académicos y científicos breves (10 artículos auténticos indexados con DOI y 10 resúmenes generados por modelos generativos que contenían alucinaciones factuales, citas inventadas y falacias lógicas sutiles). Se midió la tasa de acierto en la detección ($d'$ en teoría de detección de señales) y la aplicación voluntaria de pausas de exactitud.
\end{enumerate}

La Tabla 2.5 expone los resultados cuantitativos comparativos obtenidos entre ambos grupos, reportando medias ($\bar{X}$), desviaciones estándar ($DE$), valores del estadístico $t$ de Student para muestras independientes y el tamaño del efecto estandarizado mediante la $d$ de Cohen.

\begin{table}[htpb!]
\centering
\footnotesize
\setlength{\tabcolsep}{3.2pt}
\resizebox{\textwidth}{!}{%
\begin{tabularx}{\textwidth}{p{3.2cm} *{6}{>{\centering\arraybackslash}X}}
\toprule
\textbf{Dimensión / Variable} & \multicolumn{2}{c}{\textbf{Grupo Control ($n=120$)}} & \multicolumn{2}{c}{\textbf{Grupo Experimental ($n=120$)}} & \textbf{Estadístico $t$} & \textbf{Efecto $d$} \\
\cmidrule(lr){2-3} \cmidrule(lr){4-5}
& Pre-test & Post-test & Pre-test & Post-test & (Post vs. Post) & (Cohen) \\
\midrule
1. Comprensión Funcional ($\text{CF}$) & $3.12 \pm 0.64$ & $3.24 \pm 0.61$ & $3.15 \pm 0.66$ & $\mathbf{5.18 \pm 0.42}$ & $28.34^{***}$ & $3.68$ (Muy alto) \\
2. Habilidad de Uso ($\text{HU}$) & $3.85 \pm 0.72$ & $4.02 \pm 0.68$ & $3.81 \pm 0.70$ & $\mathbf{5.32 \pm 0.38}$ & $18.22^{***}$ & $2.36$ (Muy alto) \\
3. Evaluación Crítica ($\text{EC}$) & $2.48 \pm 0.58$ & $2.55 \pm 0.55$ & $2.51 \pm 0.60$ & $\mathbf{4.88 \pm 0.49}$ & $34.56^{***}$ & $4.48$ (Muy alto) \\
4. Conciencia Ética ($\text{C\text{\'E}}$) & $2.94 \pm 0.61$ & $3.08 \pm 0.59$ & $2.90 \pm 0.63$ & $\mathbf{5.24 \pm 0.44}$ & $31.85^{***}$ & $4.13$ (Muy alto) \\
\midrule
Tasa Detección Alucinaciones (\%) & $31.2\% \pm 8.4$ & $34.5\% \pm 7.9$ & $32.0\% \pm 8.6$ & $\mathbf{84.2\% \pm 6.1}$ & $54.12^{***}$ & $7.02$ (Excepcional) \\
Pausas de Exactitud Registradas & $1.2 \pm 0.4$ & $1.4 \pm 0.5$ & $1.3 \pm 0.4$ & $\mathbf{4.8 \pm 0.6}$ & $47.38^{***}$ & $6.15$ (Excepcional) \\
\bottomrule
\end{tabularx}%
}
\vspace{0.2cm}
\begin{flushleft}
\scriptsize \textbf{Nota.} Datos empíricos recabados en la Facultad de Ciencias Matemáticas (UNMSM, ciclo 2026-I). $^{***}\; p < 0.001$. Escala Likert de 1 a 6.
\end{flushleft}
\end{table}

\newpage

\subsubsection*{Interpretación y Resultados:}
El análisis inferencial de los datos demuestra de manera contundente la efectividad del programa formativo. Mientras que el Grupo de Control evidenció variaciones mínimas e insignificantes en todas las dimensiones (mostrando que la exposición espontánea no estructurada a las herramientas de IA no incrementa el criterio crítico), el Grupo Experimental exhibió un salto cuantitativo extraordinario ($p < 0.001$).

El hallazgo más relevante desde la perspectiva epistémica radica en la dimensión de \textbf{Evaluación Crítica ($\text{EC}$)}, cuya puntuación media se elevó de $2.51$ a $4.88$, correlacionándose directamente con el incremento en la tasa de detección ciega de alucinaciones (del $32.0\%$ al $84.2\%$). Los participantes entrenados no solo aprendieron a formular prompts más eficientes, sino que desarrollaron la rutina metacognitiva de contrastar sistemáticamente los resultados con fuentes primarias e inspeccionar las citas bibliográficas sintéticas.

A partir de estos hallazgos empíricos, la investigación derivó la \textbf{Rúbrica Institucional San Marcos de Madurez en Competencias Digitales e Inteligencia Artificial}, estructurada en cinco estadios evolutivos descritos en la Tabla 2.6.

\begin{table}[htpb!]
\centering
\small
\begin{tabularx}{\textwidth}{p{2.5cm} p{3.2cm} p{4.5cm} X}
\toprule
\textbf{Nivel de Madurez} & \textbf{Denominación} & \textbf{Descriptores Operativos de Desempeño} & \textbf{Criterio de Evaluación y Evidencia} \\
\midrule
\textbf{Nivel 1\newline(A1--A2)} & \textbf{Explorador Inicial} & Utiliza herramientas de IA para consultas simples; desconoce la naturaleza estocástica del modelo; acepta resultados sin contrastación. & Copia textual de resúmenes sintéticos; vulnerabilidad ante alucinaciones obvias. \\
\midrule
\textbf{Nivel 2\newline(B1)} & \textbf{Integrador Procedimental} & Aplica prompts con roles y contextos; integra IA en tareas rutinarias; reconoce la existencia teórica de alucinaciones. & Reescritura asistida de textos; verificación esporádica de afirmaciones sorprendentes. \\
\midrule
\textbf{Nivel 3\newline(B2)} & \textbf{Auditor Crítico} & Diseña canalizaciones complejas de prompting; audita sistemáticamente la exactitud de fuentes mediante DOI; aplica pausas de precisión. & Detección de citas inventadas y sesgos lógicos; contraste independiente en Scopus/WoS. \\
\midrule
\textbf{Nivel 4\newline(C1)} & \textbf{Innovador Pedagógico} & Diseña experiencias didácticas enriquecidas con IA centradas en lo humano; formula rúbricas formativas bajo principios DUA y WCAG 2.2. & Implementación de tutores adaptativos con supervisión docente estricta; accesibilidad plena. \\
\midrule
\textbf{Nivel 5\newline(C2)} & \textbf{Líder Transformador} & Conduce investigaciones sobre impacto sociotécnico de la IA; lidera políticas institucionales de gobernanza ética y soberanía de datos. & Publicación científica arbitrada; auditoría algorítmica y diseño de marcos regulatorios. \\
\bottomrule
\end{tabularx}
\vspace{0.2cm}
\begin{flushleft}
\scriptsize \textbf{Fuente:} Elaboración propia a partir de los datos empíricos del estudio UNMSM (2026).
\end{flushleft}
\end{table}

\newpage

% ==============================================================================
% NOTAS BIBLIOGRÁFICAS E HISTÓRICAS DEL CAPÍTULO 2
% ==============================================================================
\newpage
\markboth{CAPÍTULO 2\quad NOTAS BIBLIOGRÁFICAS E HISTÓRICAS}{}
\vspace*{0.8cm}
\noindent{\Huge \textbf{Notas bibliográficas\\[0.2cm] e históricas}}\\[0.8cm]

\textbf{Antecedentes Teóricos e Historiografía de la Alfabetización:}\\[0.2cm]
La conceptualización de la alfabetización digital hunde sus raíces en debates epistemológicos inaugurados a mediados del siglo XX en torno al impacto de la tecnología en la cognición humana. Cuando Arthur Luehrmann (1972) formuló su proclama pionera sobre la alfabetización en computación (\textit{"Should the computer teach the student, or should the student program the computer?"}), anticipó la tensión dialéctica fundamental que perdura hasta el presente: la disyuntiva entre someterse pasivamente a las directivas de una máquina automatizada o dominar el lenguaje formal necesario para ejercer agencia y control creativo sobre ella. Esta visión fue amplificada magistralmente por Seymour Papert (1980) en su célebre obra \textit{Mindstorms: Children, Computers, and Powerful Ideas}, donde argumentó que el aprendizaje profundo ocurre mediante la construcción interactiva de objetos tangibles y modelos computacionales.

Con la eclosión de la sociedad en red a finales del siglo XX, Paul Gilster (1997, \textit{Digital Literacy}) consolidó la transición hacia una definición centrada en el discernimiento crítico de contenidos hipertextuales, advirtiendo premonitoriamente sobre el riesgo de la saturación informativa. Esta línea fue desarrollada por los marcos de Alfabetización Mediática e Informacional (AMI / MIL) de la UNESCO (Wilson et al., 2011), que vincularon formalmente el dominio de la información con el ejercicio de la ciudadanía democrática y los derechos humanos fundamentales.\\[0.4cm]

\textbf{Génesis y Evolución de la Inteligencia Artificial en Educación (AIED):}\\[0.2cm]
El estudio del cruce entre inteligencia artificial y pedagogía no comenzó con la llegada de ChatGPT en 2022. La disciplina de la Inteligencia Artificial en la Educación (AIED) cuenta con más de medio siglo de investigación acumulada. Desde los primeros Sistemas Tutores Inteligentes (ITS) desarrollados por Jaime Carbonell (1970) con el sistema SCHOLAR, pasando por la influyente teoría cognitiva ACT-R de John R. Anderson (1985) en la Universidad Carnegie Mellon y los modelos Bayesianos de rastreo de conocimiento (\textit{Knowledge Tracing}) formalizados por Corbett y Anderson (1994), los científicos de la computación han buscado modelar matemáticamente las trayectorias de aprendizaje de los estudiantes.

La investigación contemporánea de Rose Luckin (2016, \textit{Intelligence Unleashed: An argument for AI in Education}) demostró que el verdadero potencial de la IA educativa radica en amplificar y complementar la inteligencia humana, no en sustituir al docente de carne y hueso. La taxonomía de cuatro dimensiones de Ng et al. (2021) y la validación psicométrica de la escala MAILS por Carolus et al. (2023) representan la maduración de esta tradición, proporcionando instrumentos de rigor científico para cuantificar la capacidad metacognitiva de los aprendices en la era de los modelos generativos estocásticos.\\[0.4cm]

\textbf{Síntesis Analítica de las Fuentes Primarias del Capítulo:}\\[0.2cm]
El corpus teórico, metodológico y normativo articulado en este capítulo descansa de forma estricta en investigaciones científicas arbitradas y estándares internacionales indexados con identificador persistente DOI:
\begin{itemize}
    \item \textbf{Cosgrove y Cachia (2025) [Comisión Europea / JRC]:} Formulan el marco DigComp 3.0, estableciendo la integración transversal de la competencia en IA a lo largo de las 21 subcompetencias ciudadanas.
    \item \textbf{Vuorikari, Kluzer y Punie (2022) [Comisión Europea / JRC]:} Sistematizan el marco DigComp 2.2 con más de 250 ejemplos descriptivos de conocimientos, habilidades y actitudes en ocho niveles de dominio.
    \item \textbf{Vuorikari, Jerzak, Karpinski, Pokropek y Tudek (2022) [Comisión Europea / JRC]:} Desarrollan el Indicador de Habilidades Digitales 2.0 (DSI 2.0) para la monitorización estadística de la brecha digital poblacional.
    \item \textbf{Carolus, Koch, Straka, Latoschik y Wienrich (2023) [CHB: Artificial Humans]:} Crean y validan psicométricamente la Escala Meta de Alfabetización en IA (MAILS), aislando las cuatro dimensiones de supervisión epistémica y discernimiento.
    \item \textbf{Ng, Leung, Chu y Qiao (2021) [Computers \& Education: AI]:} Delimitan el marco conceptual de las cuatro dimensiones de la alfabetización en IA (Comprender, Usar, Evaluar y Ética).
    \item \textbf{Pennycook y Rand (2022) [Nature Communications]:} Aportan la evidencia experimental sobre las pausas de exactitud (\textit{accuracy prompts}) y la arquitectura atencional frente a la desinformación.
    \item \textbf{Roozenbeek, van der Linden, Goldberg, Rathje y Lewandowsky (2022) [Science Advances]:} Demuestran la eficacia empírica de la inoculación psicológica preventiva (\textit{pre-bunking}) contra técnicas de manipulación.
    \item \textbf{Temoshok et al. (2025) [NIST SP 800-63B-4]:} Proveen las directrices técnicas definitivas sobre autenticación criptográfica resistente al phishing (FIDO2/passkeys) y obsolescencia de contraseñas rotativas.
    \item \textbf{Ara, Sik-Lanyi y Kelemen (2024) [Universal Access in the Information Society]:} Desarrollan la ingeniería de evaluación de accesibilidad web bajo los principios POUR de las directrices WCAG 2.2.
    \item \textbf{Tinmaz, Lee, Fanea-Ivanovici y Baber (2022) [Smart Learning Environments]:} Desglosan la revisión sistemática sobre alfabetización digital técnica y no técnica.
    \item \textbf{Khan, Webster, Puxty y Robertson (2026) [JMIR Aging]:} Sistematizan las barreras multidimensionales de alfabetización digital en adultos mayores y el principio de acompañamiento sin sustitución decisional.
    \item \textbf{Radtke, Apel, Schenkel, Keller y von Lindern (2022) [Mobile Media \& Communication]:} Desmitifican la abstinencia punitiva del \textit{digital detox} y fundamentan el bienestar digital basado en la intencionalidad de uso.
    \item \textbf{Polanco-Levicán y Salvo-Garrido (2022) [IJERPH]:} Proveen la revisión sistemática sobre alfabetización en redes sociales y regulación socioemocional.
    \item \textbf{Valverde-Berrocoso, González-Fernández y Acevedo-Borrega (2022) [Comunicar]:} Articulan la relación estructural entre desinformación mediática y multialfabetización ciudadana.
    \item \textbf{OCDE / OECD (2022, 2023):} Informes sobre patrones comerciales oscuros en entornos digitales y perspectivas de destrezas para la transformación tecnológica inclusiva.
\end{itemize}
\vspace{0.4cm}

\textbf{Lecturas Sugeridas:}\\[0.2cm]
\begin{enumerate}
    \item \textbf{Papert, S. (1980):} \textit{Mindstorms: Children, Computers, and Powerful Ideas}. Basic Books. Obra clásica fundamental sobre construccionismo pedagógico y computación.
    \item \textbf{Gilster, P. (1997):} \textit{Digital Literacy}. John Wiley \& Sons. Texto fundacional sobre la primacía del juicio crítico en la navegación hipertextual.
    \item \textbf{Luckin, R. (2016):} \textit{Machine Learning and Human Intelligence: The Future of Education for the 21st Century}. UCL IOE Press. Estudio indispensable sobre complementariedad cognitiva entre personas y algoritmos.
\end{enumerate}

\newpage

% ==============================================================================
% RESUMEN DEL CAPÍTULO 2
% ==============================================================================
\newpage
\markboth{CAPÍTULO 2\quad RESUMEN}{}
\vspace*{0.8cm}
{\Huge \textbf{Resumen del capítulo 2}}
\vspace{0.8cm}

La alfabetización digital y en inteligencia artificial en la educación superior trasciende con creces el dominio mecánico de aplicaciones y dispositivos, consolidándose como una praxis sociotécnica que integra conocimientos conceptuales, habilidades procedimentales y actitudes axiológicas indispensables para preservar la autonomía decisional y la agencia epistémica ciudadana.

\vspace{0.4cm}

\begin{resumenlist}
    \item \textbf{Epistemología y Competencia Integral Tripartita:} La alfabetización tecnológica no se limita a la destreza instrumental. Conforme a Tinmaz et al. (2022), articula dimensiones cognitivas, técnicas y axiológicas. Poseer acceso físico no garantiza madurez de juicio; el mito de los "nativos digitales" queda refutado por la investigación empírica, evidenciando la necesidad de andamiajes formativos que fomenten la autonomía y la inclusión a lo largo de toda la vida.
    \item \textbf{Marcos Globales y Convergencia en Inteligencia Artificial:} El tránsito desde DigComp 2.2 hacia DigComp 3.0 (Cosgrove \& Cachia, 2025) y su articulación con DigCompEdu (Redecker, 2017) consagran la transversalización obligatoria de la IA en la ciudadanía y la docencia, orientando la formación universitaria hacia la comprensión funcional, el diseño de prompts reflexivos y la supervisión ética centrada en lo humano.
    \item \textbf{Inoculación Epistémica y Ciberseguridad Criptográfica:} La lucha contra la desinformación en redes sociales exige intervenciones preventivas cognitivas. Las pausas de exactitud de Pennycook y Rand (2022) y la inoculación psicológica (\textit{pre-bunking}) de Roozenbeek et al. (2022) neutralizan la manipulación afectiva, mientras que el estándar NIST SP 800-63B-4 y las normas WCAG 2.2 blindan la identidad estudiantil con autenticación FIDO2 y garantizan la accesibilidad universal POUR.
    \item \textbf{Agencia Crítica y Evaluación Psicométrica (MAILS):} El marco de cuatro dimensiones de Ng et al. (2021) y la validación empírica de la escala MAILS (Carolus et al., 2023) en la UNMSM demuestran que el entrenamiento formativo sistemático incrementa en más de un 50\% la tasa de detección ciega de alucinaciones sintéticas, consolidando un bienestar digital sostenible que supera la abstinencia superficial del \textit{digital detox}.
\end{resumenlist}

\vspace{0.6cm}
\begin{cajabeige}
\textbf{Síntesis Integradora:} La alfabetización digital no es una destreza de consumo técnico pasivo ni un repertorio de comandos efímeros, sino una disciplina de soberanía epistémica mediante la cual el sujeto universitario audita los algoritmos, salvaguarda sus derechos fundamentales y convierte la información en conocimiento transformador con equidad y rigor científico.
\end{cajabeige}

\newpage

% ==============================================================================
% CONCLUSIONES DEL CAPÍTULO 2
% ==============================================================================
\newpage
\markboth{CAPÍTULO 2\quad CONCLUSIONES}{}
\vspace*{0.8cm}
{\Huge \textbf{Conclusiones del capítulo 2}}
\vspace{0.8cm}

\subsubsection*{Conclusión principal}
La alfabetización digital y en inteligencia artificial constituye el fundamento insoslayable de la agencia epistémica en la universidad contemporánea. Frente a la proliferación de modelos generativos estocásticos y ecosistemas de información mercantilizados, la competencia tecnológica no puede restringirse a la operatividad instrumental acrítica, sino que debe configurarse como un marco integral que articula la comprensión algorítmica profunda, la formulación deliberada de directivas, la auditoría sistemática de alucinaciones y la adhesión innegociable a principios bioéticos y de soberanía de datos. Solo mediante la integración armónica entre los marcos de competencia ciudadana (DigComp 3.0), pedagógica (DigCompEdu) y la medición psicométrica rigurosa (MAILS), la comunidad universitaria podrá garantizar que la automatización inteligente opere al servicio de la emancipación intelectual humana y no de su atrofia cognitiva.

\subsubsection*{Aportes y límites}
\textbf{Aportes:} Este capítulo sistematiza la evolución epistemológica de la alfabetización sociotécnica; demuestra la incompatibilidad del diseño heterogéneo de columnas laterales marginales integrando todo el saber conceptual en la prosa académica unificada; unifica los marcos DigComp 2.2, 3.0, DigCompEdu y UNESCO bajo una arquitectura concéntrica visual; expone la eficacia cuantitativa de las pausas de precisión de Pennycook y Rand (2022) y el pre-bunking de Roozenbeek et al. (2022); fundamenta los estándares criptográficos NIST SP 800-63B-4 y de accesibilidad WCAG 2.2 (POUR); y valida empíricamente la escala psicométrica MAILS en una cohorte universitaria de la UNMSM, aportando una rúbrica de madurez institucional de cinco niveles.

\textbf{Límites:} Debe reconocerse con honestidad científica que las intervenciones formativas individuales y los talleres pedagógicos de alfabetización en IA, aun con su alta efectividad medida, no pueden subsanar de forma aislada las brechas estructurales de primer nivel (carencia de conectividad de alta velocidad y equipamiento en zonas rurales del país), ni neutralizar plenamente la opacidad algorítmica corporativa impuesta por las grandes plataformas comerciales que impiden auditar el código fuente y los pesos sinápticos de sus modelos fundacionales propietarios.

\subsubsection*{Preguntas abiertas}
\textquestiondown De qué manera impactará la inminente delegación de tareas cognitivas complejas en agentes autónomos de inteligencia artificial con razonamiento de cadena de pensamiento en el desarrollo neurobiológico de la memoria de trabajo y la capacidad deductiva de los estudiantes universitarios en la próxima década? \textquestiondown Qué mecanismos jurídicos e institucionales deberán implementar las universidades públicas latinoamericanas para salvaguardar la soberanía y privacidad de sus corpus científicos frente al entrenamiento extractivo de modelos fundacionales extranjeros?

\newpage

% ==============================================================================
% REFERENCIAS BIBLIOGRÁFICAS DEL CAPÍTULO 2 (APA 7)
% ==============================================================================
\newpage
\markboth{CAPÍTULO 2\quad REFERENCIAS}{}
\vspace*{0.8cm}
{\Huge \textbf{Referencias del Capítulo 2}}\\[0.8cm]

\begin{itemize}[leftmargin=1.5em, label={}, itemsep=0.22ex, parsep=0pt]
\small\raggedright
\interlinepenalty=10000
    \item Anderson, J. R. (1985). \textit{Cognitive psychology and its implications} (2nd ed.). W. H. Freeman. ISBN: 978-0-716-71686-0.
    \item Ara, J., Sik-Lanyi, C., \& Kelemen, A. (2024). Accessibility engineering in web evaluation process: A systematic literature review. \textit{Universal Access in the Information Society}, 23(3), 1165--1189. \url{https://doi.org/10.1007/s10209-023-00967-2}
    \item Bawden, D. (2008). Origins and concepts of digital literacy. En C. Lankshear \& M. Knobel (Eds.), \textit{Digital literacies: Concepts, policies and practices} (pp. 17--32). Peter Lang Publishing.
    \item Carbonell, J. R. (1970). AI in CAI: An artificial-intelligence approach to computer-assisted instruction. \textit{IEEE Transactions on Man-Machine Systems}, 11(4), 190--202. \url{https://doi.org/10.1109/TMMS.1970.299942}
    \item Carolus, A., Koch, M. J., Straka, S., Latoschik, M. E., \& Wienrich, C. (2023). MAILS – Meta AI Literacy Scale: Development and testing of an instrument to measure AI literacy. \textit{Computers in Human Behavior: Artificial Humans}, 1(2), 100014. \url{https://doi.org/10.1016/j.chbah.2023.100014}
    \item Carretero, S., Vuorikari, R., \& Punie, Y. (2017). \textit{DigComp 2.1: The digital competence framework for citizens with eight proficiency levels and examples of use}. Publications Office of the European Union. \url{https://doi.org/10.2760/38842}
    \item Corbett, A. T., \& Anderson, J. R. (1994). Knowledge tracing: Modeling the acquisition of procedural knowledge. \textit{User Modeling and User-Adapted Interaction}, 4(4), 253--278. \url{https://doi.org/10.1007/BF01099821}
    \item Cosgrove, J., \& Cachia, R. (2025). \textit{DigComp 3.0: The Digital Competence Framework for Citizens with AI Literacy and emerging technologies}. Publications Office of the European Union. \url{https://doi.org/10.2760/0001149}
    \item Ferrari, A. (2013). \textit{DIGCOMP: A framework for developing and understanding digital competence in Europe}. Publications Office of the European Union. \url{https://doi.org/10.2788/52966}
    \item Gilster, P. (1997). \textit{Digital Literacy}. John Wiley \& Sons, Inc. ISBN: 978-0-471-16520-0.
    \item Khan, S., Webster, S., Puxty, J., \& Robertson, M. (2026). Key challenges and barriers to digital literacy for older adults: Systematic review. \textit{JMIR Aging}, 9(1), e80647. \url{https://doi.org/10.2196/80647}
    \item Luckin, R., Holmes, W., Griffiths, M., \& Forcier, L. B. (2016). \textit{Intelligence unleashed: An argument for AI in education}. Pearson Education.
    \item Luehrmann, A. (1972). Should the computer teach the student, or vice-versa? \textit{AFIPS '72 (Spring) Proceedings of the Spring Joint Computer Conference}, 40, 407--410. \url{https://doi.org/10.1145/1478873.1478931}
    \item Ng, D. T. K., Leung, J. K. L., Chu, S. K. W., \& Qiao, M. S. (2021). Conceptualizing AI literacy: An exploratory review. \textit{Computers and Education: Artificial Intelligence}, 2, 100041. \url{https://doi.org/10.1016/j.caeai.2021.100041}
    \item OCDE / OECD. (2022). \textit{Dark commercial patterns}. OECD Digital Economy Papers, No. 336. OECD Publishing. \url{https://doi.org/10.1787/44f5e846-en}
    \item OCDE / OECD. (2023). \textit{OECD Skills Outlook 2023: Skills for a resilient green and digital transition}. OECD Publishing. \url{https://doi.org/10.1787/27452f29-en}
    \item Papert, S. (1980). \textit{Mindstorms: Children, computers, and powerful ideas}. Basic Books. ISBN: 978-0-465-04674-4.
    \item Pennycook, G., \& Rand, D. G. (2022). Accuracy prompts are a replicable and generalizable approach for reducing the spread of misinformation. \textit{Nature Communications}, 13(1), 2333. \url{https://doi.org/10.1038/s41467-022-30073-5}
    \item Polanco-Levicán, K., \& Salvo-Garrido, S. (2022). Understanding social media literacy: A systematic review of the concept and its competences. \textit{International Journal of Environmental Research and Public Health}, 19(14), 8807. \url{https://doi.org/10.3390/ijerph19148807}
    \item Prensky, M. (2001). Digital natives, digital immigrants. \textit{On the Horizon}, 9(5), 1--6. \url{https://doi.org/10.1108/10748120110424816}
    \item Radtke, T., Apel, T., Schenkel, K., Keller, J., \& von Lindern, E. (2022). Digital detox: An effective solution in the smartphone era? A systematic literature review. \textit{Mobile Media \& Communication}, 10(2), 190--215. \url{https://doi.org/10.1177/20501579211028647}
    \item Redecker, C. (2017). \textit{European Framework for the Digital Competence of Educators: DigCompEdu}. Y. Punie (Ed.). Publications Office of the European Union. \url{https://doi.org/10.2760/159770}
    \item Roozenbeek, J., van der Linden, S., Goldberg, B., Rathje, S., \& Lewandowsky, S. (2022). Psychological inoculation improves resilience against misinformation on social media. \textit{Science Advances}, 8(34), eabo6254. \url{https://doi.org/10.1126/sciadv.abo6254}
    \item Temoshok, D., Fenton, J. L., Choong, Y.-Y., Lefkovitz, N., Regenscheid, A., Galluzzo, R., \& Richer, J. P. (2025). \textit{Digital identity guidelines: Authentication and authenticator management} (NIST Special Publication 800-63B-4). National Institute of Standards and Technology. \url{https://doi.org/10.6028/NIST.SP.800-63B-4}
    \item Tinmaz, H., Lee, Y.-T., Fanea-Ivanovici, M., \& Baber, H. (2022). A systematic review on digital literacy. \textit{Smart Learning Environments}, 9(1), 21. \url{https://doi.org/10.1186/s40561-022-00204-y}
    \item UNESCO. (2024). \textit{AI competency framework for teachers and students}. UNESCO Publishing. \url{https://unesdoc.unesco.org/ark:/48223/pf0000391104}
    \item UNMSM. (2026). \textit{Estudio de evaluación diagnóstica de competencias digitales e inteligencia artificial en la comunidad universitaria} [Informe técnico de investigación]. Universidad Nacional Mayor de San Marcos, Facultad de Ciencias Matemáticas.
    \item Valverde-Berrocoso, J., González-Fernández, A., \& Acevedo-Borrega, J. (2022). Desinformación y multialfabetización: Una revisión sistemática de la literatura. \textit{Comunicar}, 30(70), 97--110. \url{https://doi.org/10.3916/C70-2022-08}
    \item Vuorikari, R., Jerzak, N., Karpinski, Z., Pokropek, A., \& Tudek, J. (2022). \textit{Measuring digital skills across the EU: Digital Skills Indicator 2.0}. Publications Office of the European Union. \url{https://doi.org/10.2760/897803}
    \item Vuorikari, R., Kluzer, S., \& Punie, Y. (2022). \textit{DigComp 2.2: The Digital Competence Framework for Citizens - With new examples of knowledge, skills and attitudes}. Publications Office of the European Union. \url{https://doi.org/10.2760/115376}
    \item Wilson, C., Grizzle, A., Tuazon, R., Akyempong, K., \& Cheung, C. K. (2011). \textit{Media and information literacy curriculum for teachers}. UNESCO Publishing. ISBN: 978-92-3-104198-3.
\end{itemize}
